% BHU PYQ -- Geography: Human Geography (2024-25 NEP)
\documentclass[11pt,a4paper]{article}
\usepackage[a4paper,top=2.5cm,bottom=2.5cm,left=2.8cm,right=2.8cm,headheight=14pt]{geometry}
\usepackage{amsmath,amssymb}
\usepackage{fontspec}
\setmainfont{Kohinoor Devanagari}
\usepackage{microtype,fancyhdr,enumitem,hyperref,lastpage,parskip}

\hypersetup{colorlinks=true,urlcolor=black,linkcolor=black,pdftitle={GGRMJ21: Human Geography -- BHU NEP 2024-25}}

\pagestyle{fancy}\fancyhf{}
\renewcommand{\headrulewidth}{0.4pt}\renewcommand{\footrulewidth}{0.4pt}
\fancyhead[L]{\small \textit{Banaras Hindu University}}
\fancyhead[R]{\small \textit{GGRMJ21: Human Geography (2024-25)}}
\fancyfoot[C]{\small Page \thepage\ of \pageref{LastPage}}

\newlist{parts}{enumerate}{2}
\setlist[parts,1]{label=(\alph*),leftmargin=*,itemsep=4pt,topsep=3pt}
\setlist[parts,2]{label=(\roman*),leftmargin=*,itemsep=2pt,topsep=2pt}

\newcommand{\pts}[1]{\hfill \textit{[#1]}}

\begin{document}

\begin{center}
  {\large\bfseries BANARAS HINDU UNIVERSITY}\\[2pt]
  {\normalsize B. A. / B. Sc. (Hons.) Semester II Examination, 2024-25 (NEP)}\\[6pt]
  \rule{0.8\linewidth}{0.5pt}\\[6pt]
  {\large\bfseries GEOGRAPHY}\\[4pt]
  {\large\bfseries Paper Code: GGRMJ21 / GGRMN21 : Human Geography}\\[4pt]
  {\normalsize Time: 2:30 Hours \hfill Full Marks: 70}\\[2pt]
  {\small REF NO. CE/Conf./D-II/088 \& 089}
\end{center}
\rule{\linewidth}{0.4pt}
\smallskip

\noindent \textit{Instructions: Answer five questions in all. All questions carry equal marks (14 marks each). Question No.\ 1 is compulsory. Answer each part of Question No.\ 1 in approximately 200 words and remaining questions in approximately 1000 words. Illustrate your answer with suitable maps and diagrams. Write your Roll No.\ at the top immediately on receipt of this question paper.}

\smallskip
\noindent \textit{कुल पाँच प्रश्नों के उत्तर दीजिए। सभी प्रश्नों के अंक समान हैं। प्रश्न सं० 1 अनिवार्य है। प्रश्न संख्या 1 के प्रत्येक भाग का उत्तर लगभग 200 शब्दों में एवं शेष प्रश्नों का उत्तर लगभग 1000 शब्दों में दीजिए। अपने उत्तर की पुष्टि उपयुक्त मानचित्रों तथा आरेखों द्वारा कीजिए।}

\bigskip

\begin{enumerate}[label=\textbf{\arabic*.},leftmargin=*,itemsep=14pt]

  \item Write short notes on the following:\pts{4 $\times$ 3.5 = 14}
  
  निम्नलिखित पर संक्षिप्त टिप्पणियाँ लिखिए:
  \begin{parts}
    \item Evolution of races / प्रजातियों का विकास
    \item Habitat of Naga tribe / नागा जनजाति का निवास क्षेत्र
    \item Trend of world population growth / विश्व जनसंख्या वृद्धि की प्रवृत्ति
    \item Functional classification of urban settlement / नगरीय बस्तियों का कार्यात्मक वर्गीकरण
  \end{parts}

  \item Discuss the nature and scope of human geography.\pts{14}
  
  मानव भूगोल की प्रकृति एवं विषय-क्षेत्र की विवेचना कीजिए।

  \item Discuss the salient aspects of `Determinism' and `Possibilism'.\pts{14}
  
  `नियतिवाद' एवं `संभववाद' के प्रमुख पक्षों की चर्चा कीजिए।

  \item Describe the characteristics and distribution of the major races of the world.\pts{14}
  
  विश्व की प्रमुख प्रजातियों की विशेषताओं एवं वितरण का वर्णन कीजिए।

  \item Give an account of the habitat, economy and society of the Eskimos.\pts{14}
  
  एस्किमो के निवास क्षेत्र, अर्थव्यवस्था एवं समाज का विवरण प्रस्तुत कीजिए।

  \item Define Migration and describe its major types with suitable examples.\pts{14}
  
  प्रवास को परिभाषित कीजिए एवं उपयुक्त उदाहरणों के साथ इसके मुख्य प्रकारों का वर्णन कीजिए।

  \item Discuss in detail the characteristics and types of rural settlements.\pts{14}
  
  ग्रामीण बस्तियों की विशेषताओं एवं प्रकारों की विस्तार से चर्चा कीजिए।

\end{enumerate}

\end{document}
